% !TEX program = lualatex
% !TeX encoding = UTF-8
\PassOptionsToPackage{unicode}{hyperref}
\PassOptionsToPackage{bookmarksnumbered,bookmarksopen}{hyperref}
\documentclass[11pt,a4paper]{article}

% ============================================================
% 한국어 설정 (LuaLaTeX + luatexja)
% ============================================================
\usepackage{luatexja}
\usepackage{luatexja-fontspec}

% 한국어 고딕체 (sans) – Noto Sans KR (Windows/Mac/Linux)
\setsansjfont{Noto Sans KR}[
UprightFont = * Regular,
BoldFont    = * Bold,
Script      = Hangul,
Language    = Korean
]

% 한국어 명조체 (serif) – Noto Serif KR
\IfFontExistsTF{Noto Serif KR}{
	\setmainjfont{Noto Serif KR}[
	UprightFont = * Regular,
	BoldFont    = * Bold,
	Script      = Hangul,
	Language    = Korean
	]
}{
	% fallback: Noto Sans KR (만약 Serif가 없으면)
	\setmainjfont{Noto Sans KR}[
	UprightFont = * Regular,
	BoldFont    = * Bold,
	Script      = Hangul,
	Language    = Korean
	]
}

% 고정폭 폰트 – 라틴 문자는 Latin Modern Mono, 한글은 Noto Sans KR
\setmonojfont{Noto Sans KR}[
UprightFont = * Regular,
BoldFont    = * Bold,
Script      = Hangul,
Language    = Korean
]

% 라틴 문자 폰트 (영문)
\setmainfont{Times New Roman}[Ligatures=TeX]
\setsansfont{Arial}[Ligatures=TeX]
\setmonofont{Latin Modern Mono}[
WordSpace={1,0,0},
HyphenChar=None,
PunctuationSpace=WordSpace
]

% ============================================================
% 기타 패키지 (원본과 동일)
% ============================================================
\usepackage{amsmath,amssymb,amsthm,mathtools,bm}
\usepackage{geometry}
\usepackage{enumitem}
\usepackage{siunitx}
\usepackage{booktabs}
\usepackage{array}
\usepackage{slashed}
\usepackage{graphicx}
\usepackage{caption}
\usepackage{subcaption}
\usepackage{braket}
\usepackage{tensor}
\usepackage{makecell}
\usepackage[most]{tcolorbox}
\usepackage{pgfplots}
\usepackage{float}
\usepackage{tikz}
\usepackage{multicol}
\usepackage{lipsum}
\usepackage{microtype}
\usepackage{hyperref}

\pgfplotsset{compat=1.18}
\usetikzlibrary{arrows.meta,positioning,shapes.geometric,fit,calc,
	decorations.pathreplacing,patterns,quotes,angles,
	shadows,decorations.markings}

\geometry{margin=1in}
\setlength{\emergencystretch}{3em}
\sloppy

% 정리 환경 (한국어)
\newtheorem{theorem}{정리}
\newtheorem{lemma}{보조정리}
\newtheorem{axiom}{공리}
\newtheorem{remark}{주석}
\newtheorem{proposition}{명제}
\newtheorem{definition}{정의}
\newtheorem{assumption}{가정}
\newtheorem{corollary}{따름정리}

% PDF 메타데이터
\hypersetup{
	pdftitle={부록 Climat01: YUCT의 기후 조정학 — 기후 변동의 보편적 스케일링, 대기 열용량, 해양 조정 및 임계 전이 예측},
	pdfauthor={Alexey V. Yakushev},
	pdfsubject={조정 효율 (K_eff), 기후 변동성, 태양 활동, 해양 열량, 빙하 질량 수지, 화산 강제력, 맹목적 예측, 임계 전이},
	pdfkeywords={YUCT, 조정 효율, 기후 변동성, 태양 활동, 해양 열량, 빙하 질량 수지, 화산 강제력, 맹목적 예측, 임계 전이},
	breaklinks=true,
	pdfencoding=auto,
	bookmarksnumbered=true,
	bookmarksopen=true,
	bookmarksopenlevel=1
}

% 매크로 (원본과 동일)
\newcommand{\Keff}{K_{\mathrm{eff}}}
\newcommand{\kappac}{\kappa_c}
\newcommand{\be}{\beta}
\newcommand{\ga}{\gamma}
\newcommand{\lun}{\mathcal{L}}

\begin{document}
	
	\title{부록 Climat01: YUCT의 기후 조정학\\
		\large --- 기후 변동의 보편적 스케일링, 대기 열용량, 해양 조정 및 임계 전이 예측 ---}
	\author{Alexey V. Yakushev}
	\date{2026년 3월}
	
	\maketitle
	
	\begin{abstract}
		\noindent
		본 부록은 야쿠셰프 통합 조정 이론(YUCT)을 지구 기후 시스템에 포괄적으로 적용한 것이다. 보편적 오차 스케일링 법칙 $\varepsilon = \kappa_c \alpha K_{\mathrm{eff}}^{-\beta}$ ($\beta = 0.67$)이 전 지구 기온, 해양 열량, 빙상 질량 손실의 수십 년 규모 변동성을 지배함을 보인다. 태양 활동(볼프 수)은 기후 시스템의 조정 효율 $K_{\mathrm{eff}}$을 변조하는 외부 지표로 작용한다. 30년 이동 창을 사용하면 $\ln \sigma_T$의 $\ln W$에 대한 회귀 기울기는 $-0.70\pm0.11$ ($R^2=0.62$)로 이론 지수와 매우 잘 일치한다. 동일한 스케일링이 해수면 온도(SST)와 해양 열량에서도 나타나며, 지역적 차이는 해양 관성과 달의 교점 주기로 설명된다. 빙권에서는 비선형 융해 과정을 반영하여 지수가 $\approx 0.5$가 된다. 화산 강제력의 새로운 모델을 도입하여, 분화 확률을 동일한 지수 $\beta=0.67$로 지자기 활동과 연결하고, 에어로졸 광학 두께를 변동성 방정식에 지수 감쇠 항으로 포함한다. 통합 모델을 사용하여 2050년까지의 지표 기온 편차에 대한 지역별 맹목적 예측을 수행하고, 2045년경 기후 시스템이 임계閾値 ($K_{\mathrm{eff}} \to K_{\mathrm{crit}}$)에 접근함을 논의한다. 본 틀은 기후 동역학을 경제 및 지구물리학적 현상과 단일한 조정 원리로 통합하여, 향후 수십 년간 검증 가능한 예측을 제공한다.
	\end{abstract}
	
	\tableofcontents
	
	\section*{공리 0: 조정의 존재론적 우선성}
	\addcontentsline{toc}{section}{공리 0: 조정의 존재론적 우선성}
	
	모든 이론은 정태(靜態)를 기술한다. 그러나 조정이야말로 정태를 가능하게 하는 과정이다.
	어떠한 명제, 모델, 법칙도 사전(辭典)의 활성화된 요소로서만 존재한다.
	사전과 색인은 그것들이 기술하는 정태로부터 도출될 수 없으며, 존재론적으로 선행한다.
	
	YPSDC 프로토콜은 "여러 이론 중 하나"가 아니다.
	그것은 자신을 정식화하는 이론을 포함한, 모든 이론의 기저에 있는 프로토콜이다.
	이를 거부하는 것은 정식화 능력 자체를 거부하지 않고는 불가능하다.
	
	그런 의미에서, 조정은 다른 이론들보다 "더 강력한" 것이 아니라 **더 근원적**이다.
	정태는 생성하지 않으며, 정태가 생성된다.
	
	이 공리는 YUCT의 근본적인 존재론적 기초로 채택된다.
	
	\section{서론}
	
	지구 기후 시스템은 고도로 조정된 복잡계이다. 대기, 해양, 빙권, 생물권은 끊임없이 에너지, 운동량, 물질을 교환하며, 안정된 홀로세 상태를 유지하려면 높은 수준의 조정이 필요하다. 야쿠셰프 통합 조정 이론(YUCT)의 틀 안에서, 모든 조정계는 삼원소 $D + I \cdot R$로 기술된다:
	
	\begin{itemize}
		\item $D$ (사전) – 유체 역학, 복사 전달, 열역학의 물리 법칙; 선험적으로 분배됨.
		\item $I$ (정보) – 현재 상태 변수: 기온, 기압, 염분, 빙하 면적.
		\item $R$ (공명) – 광대한 지역을 동기화하는 결맞은 변동 패턴 (ENSO, NAO, AMO).
	\end{itemize}
	
	기후 시스템의 조정 효율 $\Keff$은 섭동 후 평형으로 돌아가는 속도를 결정한다. 보편적 YUCT 법칙에 따라, 상대 오차(변동의 진폭)는 다음과 같이 스케일링된다:
	\[
	\epsilon = \kappac \alpha \Keff^{-\be}, \qquad \be = 0.67 \pm 0.03,\ \kappac = 0.33,
	\]
	여기서 $\epsilon$은 기후 변동성의 정규화된 진폭(예: 기온 편차의 표준 편차)으로 해석될 수 있고, $\alpha$는 시스템 특유의 상수이다.
	
	본 부록에서는 태양 활동(볼프 수 $W$)이 외부 조정 지표로 작용하여, 이 법칙이 전 지구 기온에 성립함을 보인다. 얻어진 지수는 $\be = 0.67$과 일치하며, 이전에 경제학(부록 F)에서 발견된 지수와 일치한다. 이는 YUCT의 보편성에 대한 강력한 학제적 확증을 제공한다. 나아가 해양 성분, 빙권, 대기 온도의 섹터 구조로 분석을 확장하여, 달의 조석 및 중력 의존성(부록 P)과의 연결을 확립한다.
	
	\section{이론적 틀}
	
	\subsection{기후 시스템의 조정 효율}
	
	대기는 열에너지를 운동 에너지로 변환하는 열기관으로 작동한다. 그 효율은
	\[
	\eta = \frac{T_h - T_c}{T_h}
	\]
	로 표현된다. 여기서 $T_h$는 열 입력 온도(열대 지표), $T_c$는 열 방출 온도(극지 상부 대류권)이다. YUCT에서 대기 조정 효율은 $\eta$와 다음과 같이 연결된다:
	\[
	\Keff^{\text{atm}} = K_0 \left( \frac{\eta}{\eta_0} \right)^{\be}.
	\]
	
	대기의 열에너지 저장 능력(유효 열용량)은 $\Keff$에 따라 다음과 같이 스케일링된다:
	\[
	C_{\text{heat}} = C_0 \cdot \Keff^{\ga}, \qquad \ga \approx 0.3 \pm 0.05,
	\]
	여기서 $\ga$는 부록 P에서 온도와 중력 상수를 연결하는 지수와 동일하다.
	
	\subsection{기온 변동성의 스케일링}
	
	보편 법칙으로부터, 상대 변동 진폭은 $\Keff^{\be}$에 반비례한다. 전 지구 기온에 대해:
	\[
	\frac{\sigma_T}{\langle T \rangle} \propto \Keff^{-\be}.
	\]
	$\Keff$이 태양 활동 $\langle W \rangle$ (특성 시간으로 평균)에 비례한다고 가정하면, 검증 가능한 관계를 얻는다:
	\[
	\sigma_T \propto \langle W \rangle^{-\be}.
	\]
	
	\subsection{인위적 강제력의 조정 모델}
	
	현대 기후 변화는 자연적 요인(태양 활동, 화산)과 인위적 요인(온실 기체, 에어로졸) 모두에 의해 구동된다. YUCT에서, 외부 강제력은 조정 효율에 영향을 미친다. 인위적 영향을 고려하기 위해, 복사 강제력 $\Delta F_{\mathrm{anth}}(t)$에 비례하는 항을 $K_{\mathrm{eff}}$의 진화 방정식에 추가한다:
	\[
	\frac{\partial K}{\partial t} = r K \left(1 - \frac{K}{K_{\mathrm{opt}}}\right) - \mu \kappa_c \alpha K^{1-\beta} 
	+ D\nabla^2 K + \gamma_{\mathrm{anth}} \frac{\Delta F_{\mathrm{anth}}(t)}{\Delta F_0} K^{1-\beta} + \xi(t),
	\]
	여기서 $\gamma_{\mathrm{anth}}$는 인위적 강제력에 대한 조정 효율의 민감도(역사적 데이터로부터 교정)이다. 본 계산에서는 IPCC AR6의 SSP2-4.5 및 SSP5-8.5 시나리오로부터 $\Delta F_{\mathrm{anth}}(t)$를 사용한다. 인위적 강제력의 $K_{\mathrm{eff}}$ 감소에 대한 기여는 태양 활동의 효과와 유사하여, 수십 년 진동과 장기 추세를 모두 모델링할 수 있다.
	
	\subsection{조정 과정의 결과로서의 화산 활동}
	
	대규모 화산 분화는 성층권 에어로졸을 통해 기후를 변조한다. 그러나 그 통계는 순수히 무작위적이지 않다: 역사적 및 고기후 재구성은 강력한 분화 빈도와 태양 활동 및 지자기 교란(예: 11년 주기 및 그 조화)의 위상 사이의 상관관계를 보여준다. YUCT에서 이는 $\Keff$의 변화가 대기와 해양뿐만 아니라 지구 동역학적 과정에도 영향을 미치기 때문으로 설명된다.
	
	\subsubsection{분화 방아쇠 지표로서의 지자기 지수}
	
	태양 활동을 고체 지구 과정과 연결하는 통합 지표는 **지자기 지수 $Kp$** (또는 그 파생, 예: $Ap$)이다. 지자기 교란 증가기(태양 극대기)는 특히 조석 응력에 민감한 지역에서 더 많은 대규모 분화와 상관된다.
	
	보편적 오차 스케일링 법칙으로부터, 지자기 변동의 상대 진폭 $\delta Kp / \langle Kp \rangle$는 $\Keff^{-\beta}$로 스케일링된다. 단위 시간당 분화 확률은 다음과 같이 쓸 수 있다:
	\[
	P_{\text{eruption}}(t) = P_0 \cdot \left( \frac{Kp(t)}{Kp_0} \right)^{\beta} \cdot f(\text{위도}, \text{지구조}),
	\]
	여기서 $P_0$는 기준 빈도, $f$는 지질학적 환경을 고려한 지역 인자이다. 지수가 양수인 이유는 조정 효율 상승(태양 활동 경유)이 방아쇠 확률을 높이기 때문이다.
	
	\subsubsection{결합 메커니즘}
	
	영향의 주요 경로는 다음과 같다:
	\begin{itemize}
		\item \textbf{조석 응력}: 태양 활동에 의해 $G_{\text{eff}}(T)$ (부록 P)를 통해 변조된 중력 퍼텐셜 변화가 지각 및 맨틀 내 압력 응력을 변화시켜, 이미 "준비된" 마그마계의 분화 방아쇠로 작용할 수 있다.
		\item \textbf{전자기 유도}: 지자기 폭풍은 지각 내에 전류를 유도하여 국부적 가열 및 마그마 점성 변화를 일으켜 분화 임계값을 낮춘다.
	\end{itemize}
	두 메커니즘 모두 동일한 근본 원인—$\Keff$의 변화—을 공유하며, 따라서 화산 활동에 대한 기여는 동일한 지수 $\beta$의 멱법칙을 따른다.
	
	\subsubsection{통계적 검증}
	
	역사적 분화 카탈로그(예: Smithsonian Global Volcanism Program)와 태양 활동 데이터(볼프 수)의 분석은 VEI $\ge 4$의 분화에 대해 11년 주기와의 상관관계가 $p < 0.01$로 유의함을 보여준다. 30년 창(기온 시계열과 동일하게)으로 평균하면, $\ln P_{\text{eruption}}$의 $\ln W$에 대한 회귀 기울기는 $0.65 \pm 0.15$로, $\beta = 0.67$과 일치한다.
	
	\subsubsection{에어로졸 부하 모델}
	
	분화 후 성층권 에어로졸 광학 두께는 지수 감쇠로 근사된다:
	\[
	\tau_{\text{volc}}(t) = \sum_{i} A_i \cdot \exp\!\left(-\frac{t-t_i}{\tau_{\text{dec}}}\right) \cdot g(\varphi_i, s_i),
	\]
	여기서:
	\begin{itemize}
		\item $A_i$ – 초기 에어로졸 부하 (방출된 $SO_2$ 질량에 비례);
		\item $t_i$ – 분화 시각;
		\item $\tau_{\text{dec}} \approx 1$ 년 – 성층권 에어로졸의 특성 수명;
		\item $g(\varphi_i, s_i)$ – 위도 $\varphi_i$와 계절 $s_i$를 고려한 인자 (적도 분화는 전 지구적으로 확산, 고위도 분화는 주로 자 반구에 머무름).
	\end{itemize}
	
	대규모 분화(VEI $\ge 4$)의 경우 초기 부하는 역사적 데이터(예: 위성 관측 또는 연륜 연대학)로부터 추정된다. 예측 계산에서 분화는 포아송 과정으로 모델링되며, 빈도 $\nu \approx 0.3$ yr$^{-1}$ (최근 500년 기준) 및 진폭 분포는 멱법칙을 따른다.
	
	\subsubsection{기온 변동성 방정식으로의 통합}
	
	화산 강제력을 포함하면, 전 지구 변동성 방정식은 다음과 같이 된다:
	\[
	\sigma_T(t) = \sigma_0 \left( \frac{W(t)}{W_0} \right)^{-\beta}
	\left(1 + \gamma_{\text{anth}} \frac{\Delta F_{\text{anth}}(t)}{\Delta F_0}\right)^{-\beta}
	\cdot \exp\!\bigl(-\lambda_{\text{volc}} \,\tau_{\text{volc}}(t)\bigr) \cdot \Theta_{\text{volc}}(t),
	\]
	여기서 $\lambda_{\text{volc}} > 0$는 에어로졸 부하에 대한 변동성의 민감도(역사적 데이터, 예: 1991년 피나투보 분화 후로부터 교정), $\Theta_{\text{volc}}(t)$는 분화의 확률적 성질로 인한 불확실성을 고려한다.
	
	\subsubsection{매개변수 교정}
	
	교정에는 20~21세기의 가장 큰 분화들(1991년 피나투보, 1982년 엘치촌, 1883년 크라카타우)의 데이터를 사용하였다. 최소 제곱 적합을 통해 다음을 얻었다:
	\[
	\lambda_{\text{volc}} = 0.032 \pm 0.008, \qquad \tau_{\text{dec}} = 1.1 \pm 0.2\ \text{yr}.
	\]
	$SO_2$ 질량 $Q$ (메가톤 단위)의 분화에 대한 초기 에어로졸 부하 $A_i$는 $A_i = 0.014\,Q$로 근사된다 (위성 데이터 및 빙하 코어에 기반).
	
	\subsection{YUCT를 이용한 저기압 발생 및 폭풍 강화 예측}
	\label{subsec:cyclone_prediction}
	
	YUCT 틀은 열대 저기압의 발달과 강화를 예측하는 데 사용될 수 있는 일련의 조작적 지표들을 제공한다.
	
	\subsubsection*{대기 데이터에서의 $K_{\mathrm{eff}}$ 대리}
	
	대류 영역의 조정 효율은 직접 측정할 수 없지만, 표준 기상 변수로부터 추정할 수 있다:
	\begin{itemize}
		\item \textbf{해수면 온도 (SST)}: $K_{\mathrm{eff}} \propto T_{\mathrm{SST}}^{\gamma}$, $\gamma \approx 0.3$ (부록~P). 더 따뜻한 물은 더 큰 에너지 플럭스를 제공하고 조직화된 대류를 촉진한다.
		\item \textbf{대류 가용 위치 에너지 (CAPE)}: CAPE가 높을수록 더 큰 온도 구배와 강한 연직 일관성을 의미한다.
		\item \textbf{연직 바람 시어}: 약한 시어 ($< 10$ m/s, 850–200 hPa 사이)는 소용돌이 기둥의 조정을 보존하고, 강한 시어는 이를 파괴한다.
		\item \textbf{대류권 중층의 상대 습도}: 습한 공기는 D+I·R 삼원소의 공명 인자 $R$을 강화한다.
	\end{itemize}
	합성 지표는 다음과 같이 구성될 수 있다:
	\[
	\Keff^{\mathrm{(index)}} = w_1 \left( \frac{\mathrm{SST}}{\mathrm{SST}_0} \right)^{\gamma}
	+ w_2 \left( \frac{\mathrm{CAPE}}{\mathrm{CAPE}_0} \right)^{0.67}
	+ w_3 \left( \frac{1}{1 + \mathrm{shear}/\mathrm{shear}_0} \right)
	+ w_4 \left( \frac{\mathrm{RH}}{\mathrm{RH}_0} \right),
	\]
	여기서 가중치 $w_i$는 역사적 데이터로부터 교정된다. 이 지표가 임계값 $\approx 1.84$를 초과하면, 24--48시간의 리드 타임으로 저기압 발생이 예측된다.
	
	\subsubsection*{급속 강화}
	
	위성 관측은 일부 열대 저기압이 24시간 동안 30노트 이상의 풍속 증가를 보이는 급속 강화를 겪음을 보여준다. YUCT에서 이는 자기 강화적 조정 캐스케이드에 해당한다. 일단 소용돌이가 확립되면, 그 회전이 국부적인 바람 시어를 감소시켜 $\Keff$을 더욱 증가시키고 강화를 가속한다. 이 과정은 동역학 방정식
	\[
	\frac{d\Keff}{dt} = r \Keff \left( 1 - \frac{\Keff}{K_{\max}} \right)
	- \mu \Keff^{1-\beta} + \xi(t)
	\]
	으로 기술된다. 이는 프랙털 오차 보정을 수반한 로지스틱 성장으로, 모든 조정계를 지배하는 것과 동일하다 (부록~T). 해는 가용한 해양 에너지에 의해 설정되는 $K_{\max}$에서 포화될 때까지 쌍곡선적 성장을 보인다.
	
	\subsubsection*{경험적 검증}
	
	대서양 허리케인 데이터베이스 (HURDAT2, 1851--2025)의 체계적 후향 분석은 다음을 밝혀야 한다:
	\begin{enumerate}
		\item 저기압 발생 확률은 합성 $\Keff$ 지표의 비선형 함수이며, $K_{\mathrm{crit}} \approx 1.84$ 근처에서 예리한 임계값이 존재한다.
		\item 강화율은 로지스틱 오차 방정식을 따르며, 지수 $\beta = 0.67$이 감쇠 항에 나타난다.
		\item 허리케인의 최대 잠재 강도 (MPI)는 $K_{\max}$에 비례하며, $K_{\max}$ 자체는 SST와 $\mathrm{MPI} \propto \mathrm{SST}^{\beta\gamma} \propto \mathrm{SST}^{0.20}$로 스케일링된다.
	\end{enumerate}
	이러한 예측은 반증 가능하며 공개된 데이터로 검증할 수 있다.
	
	\subsection{임계점과 전조}
	
	시스템이 임계값(임계점)에 접근함에 따라 $\Keff$이 감소하여 다음을 초래한다:
	\begin{itemize}
		\item \textbf{임계 감속}: 자기 상관 $\phi(\tau)$의 증가:
		\[
		\phi(\tau) \propto \exp\!\left(-\frac{\tau}{\tau_{\text{rel}}}\right),\quad \tau_{\text{rel}} \propto \frac{1}{\Keff}.
		\]
		\item \textbf{분산 증가}: $\sigma^2 \propto 1/\Keff$.
		\item \textbf{비대칭성 및 깜빡임}: 쌍안정 변동의 출현.
	\end{itemize}
	이러한 신호들은 급격한 기후 변화에 앞서 고기후 기록에서 관찰되며 예측에 사용될 수 있다.
	
	\subsection{소산 구조로서의 대기}
	\label{subsec:atmo_diss}
	
	지구 대기는 비평형 열역학계의 고전적인 예이다. 태양 가열은 연속적인 에너지 플럭스를 제공하여 대기를 평형에서 멀리 유지한다. 이러한 조건에서 시스템은 자체 조직화하여 개별 분자의 특성 이완 시간보다 훨씬 긴 수명을 가진 일관된 구조를 형성한다~\cite{Prigogine1977}.
	
	저기압, 허리케인, 태풍은 프리고진의 의미에서 **소산 구조**이다: 엔트로피 생성 속도 $\sigma$가 임계閾値 $\sigma_{\mathrm{crit}}$를 초과할 때 자발적으로 발생하고, 엔트로피를 환경으로 수출함으로써 질서 있는 형태를 유지하며, 에너지 공급이 중단되면 붕괴한다. YUCT에서 이 현상론은 정량적 기초를 얻는다: 허리케인의 조정 효율 $\Keff$은 보편적 스케일링 법칙을 통해 엔트로피 생성 속도와 직접 연결된다:
	\begin{equation}
		\frac{\sigma}{\sigma_0} = \left( \frac{\Keff}{K_0} \right)^{\beta d_f / 2},
		\label{eq:cyclone_sigma}
	\end{equation}
	여기서 $\beta = 2/3$, $d_f = 11/5$는 대기 조정 네트워크의 프랙털 차원이다. 지수 $\beta d_f / 2 = 0.733$이 폭풍의 에너지와 공간적 범위 사이의 멱법칙 관계를 결정한다.
	
	\subsubsection*{조정 상전이로서의 저기압 발생}
	
	열대 저기압의 탄생은 **조정 상전이**이다. 임계 해수면 온도 ($T_{\mathrm{SST}} \lesssim 26.5^{\circ}\mathrm{C}$) 이하에서는 대기 대류는 무조정적이다: 적운이 확률적으로 상승 및 소멸하며 일관된 순환을 형성하지 않는다. 해양 온도가 임계값을 초과하면 증발 속도가 증가하고, 대류 가용 위치 에너지(CAPE)가 증가하며, 시스템은 분기를 겪는다: 대규모의 고도로 조정된 소용돌이가 출현한다~\cite{Emanuel2003}.
	
	YUCT에서 이 전이는 국부 조정 효율이 임계값을 초과할 때 발생한다:
	\begin{equation}
		\Keff^{\mathrm{(atm)}} > K_{\mathrm{crit}} \quad \Longrightarrow \quad
		\text{저기압 발생}.
	\end{equation}
	임계 효율은 임의의 상수가 아니라, 대수적 루프를 통해 보편적 질서-혼돈 매개변수와 연결된다:
	\begin{equation}
		K_{\mathrm{crit}} = K_0 \left( \frac{S_{\mathrm{odd}}}{S_{\mathrm{even}}} \right)^{1/\beta}
		= K_0 \left( \frac{3}{2} \right)^{3/2} \approx 1.837\,K_0 .
	\end{equation}
	따라서 저기압은 국부적인 $K_{\mathrm{eff}}$이 기준 값을 $\sim 1.84$배 초과할 때 형성된다. 이 숫자는 YUCT의 대수적 구조에 의해 엄밀하게 결정되며, 어떤 경험적 적합에도 의존하지 않는다.
	
	\section{데이터 및 방법}
	
	\subsection{데이터 출처}
	
	\begin{itemize}
		\item \textbf{전 지구 기온}: 육지 및 해양 표면 편차 (GISTEMP, NASA) \cite{GISTEMP}. 1880년부터 2024년까지의 연간 값.
		\item \textbf{태양 활동}: 연간 볼프 수 버전 2.0 (WDC‑SILSO, 브뤼셀) \cite{SILSO}.
		\item \textbf{해양 온도}: ERSSTv5 (NOAA) 및 HadSST4 (Met Office) \cite{ERSST, HADSST}, 1850년 이후.
		\item \textbf{해양 열량 (OHC)}: IAP/CAS 데이터 \cite{IAPOHC} (0–2000 m).
		\item \textbf{빙하 및 빙상}: WGMS \cite{WGMS}, IMBIE \cite{IMBIE} 및 GRACE/GRACE-FO \cite{GRACE}에 의한 질량 수지.
		\item \textbf{섹터별 기온}: GISTEMP 띠 평균 (90°S–90°N, 위도대) \cite{GISTEMPZonal}.
		\item \textbf{화산 활동}: 전 지구 분화 카탈로그 (VEI $\ge 4$) 및 에어로졸 부하 재구축 (예: 빙하 코어 데이터) \cite{VolcanoCatalog}.
	\end{itemize}
	
	\subsection{계산 방법 및 통계적 검정}
	
	각 이동 창 크기 $L$ (10년부터 50년까지, 1년 단계)에 대해:
	\begin{enumerate}
		\item 추세 제거된 기온의 표준 편차 $\sigma_T(L,t)$와 평균 볼프 수 $\langle W(L,t) \rangle$를 각 연도 $t$에 대해 계산.
		\item $\ln \sigma_T$의 $\ln \langle W \rangle$에 대한 선형 회귀 수행.
		\item 기울기 $b(L)$, 95% 신뢰 구간 (블록 부트스트랩 10,000회, 자기 상관을 고려하여 블록 크기 $L/3$), $p$값을 기록.
	\end{enumerate}
	
	안정성 점검:
	\begin{itemize}
		\item $\ln \sigma_T$와 $\ln \langle W \rangle$에 대한 정상성 검정 (ADF, KPSS).
		\item 장기 관계에 대한 Johansen 공적분 검정.
		\item Granger 인과성 검정 (래그 1–5)으로 방향성 평가.
		\item 순열 검정 (10,000 셔플)으로 태양 활동 시계열의 무작위 재배열을 통해 $\le -0.67$의 기울기를 얻을 확률 평가.
	\end{itemize}
	
	화산 성분에 대해:
	\begin{itemize}
		\item 분화 빈도 (VEI $\ge 4$, 30년 창으로 평활화)와 볼프 수의 상관 분석.
		\item 자기 상관을 고려한 $\ln P_{\text{eruption}}$의 $\ln W$에 대한 회귀.
	\end{itemize}
	
	\subsection{강건성 평가}
	
	창 고유의 튜닝을 배제하기 위해, $b(L)$의 창 크기 의존성을 조사하였다. 기울기가 넓은 범위의 창에서 안정적이면 객관적인 규칙성을 나타낸다.
	
	\section{결과}
	
	\subsection{예비 통계 검정}
	
	ADF 검정: 두 시계열 모두 1차 차분 후 정상 ($p < 0.01$). KPSS도 차분의 정상성을 확인.  
	Johansen 검정 (랭크 1): 대각합 통계량은 20–40년 창에서 5% 수준으로 $\ln \sigma_T$와 $\ln \langle W \rangle$ 간의 공적분을 나타냄.  
	Granger 인과성: 20–40년 창에서 태양 활동에서 기온 변동성으로의 일방향 인과성 검출 ($p < 0.05$); 역방향은 비유의적.
	
	\subsection{$b(L)$의 창 크기 의존성}
	
	\begin{table}[htbp]
		\centering
		\caption{다른 창에 대한 $\ln \sigma_T$의 $\ln \langle W \rangle$에 대한 회귀 기울기}
		\label{tab:slopes}
		\begin{tabular}{@{}lcccc@{}}
			\toprule
			창 (yr) & 기울기 $b$ & 95\% CI (부트스트랩) & $p$값 & $R^2$ \\
			\midrule
			11 & –0.32 & [–0.48, –0.16] & 0.023 & 0.18 \\
			20 & –0.59 & [–0.79, –0.39] & 0.001 & 0.51 \\
			30 & –0.70 & [–0.92, –0.48] & $3\times10^{-4}$ & 0.62 \\
			40 & –0.64 & [–0.88, –0.40] & 0.002 & 0.58 \\
			\bottomrule
		\end{tabular}
	\end{table}
	
	20–40년 창에서는 기울기가 안정적이며 –0.62 … –0.72 범위에 있어 이론값 $\be = 0.67$과 완전히 일치한다. 짧은 창 (10–14년)에서는 기울기가 더 얕은데, 이는 기온 기록에 지배적인 11년 성분이 존재하지 않음에 대응한다.  
	30년 창에 대한 순열 검정은 우연히 $\le -0.67$의 기울기를 얻을 확률이 $<0.002$임을 보여, 우연의 일치라는 귀무 가설을 기각한다.
	
	\begin{figure}[htbp]
		\centering
		\begin{tikzpicture}
			\begin{axis}[
				xlabel = {창 크기 $L$ (yr)},
				ylabel = {기울기 $b$},
				grid = both,
				xmin = 10, xmax = 50,
				ymin = -0.9, ymax = -0.2,
				legend pos = north west
				]
				\addplot[thick, blue, mark=*] coordinates {
					(11, -0.32) (15, -0.48) (20, -0.59) (25, -0.66) (30, -0.70) (35, -0.68) (40, -0.64) (45, -0.60) (50, -0.55)
				};
				\addplot[thick, red, dashed] coordinates {(10, -0.67) (50, -0.67)};
				\legend{계산된 기울기, $\be = 0.67$ (참조)}
			\end{axis}
		\end{tikzpicture}
		\caption{이동 창 크기에 대한 회귀 기울기의 의존성. 넓은 창 범위 (20–40년)에서 $-\be$에 가까운 기울기를 제공한다.}
		\label{fig:slope_vs_window}
	\end{figure}
	
	\subsection{경제학과의 비교 (부록 F)}
	
	부록 F에서 GDP의 변동성 (5년 창)은 $-0.67 \pm 0.05$의 기울기를 주었다. 기후학에서는 30년 창(수십 년 진동에 대응)이 $-0.70 \pm 0.11$의 기울기를 준다. 시간 척도는 다르지만 지수는 동일하며, $\be = 0.67$의 보편성을 확증한다.
	
	\begin{table}[htbp]
		\centering
		\caption{스케일링 지수의 비교}
		\label{tab:compare}
		\begin{tabular}{@{}lccc@{}}
			\toprule
			시스템 & 창 & 물리적 의미 & 기울기 \\
			\midrule
			경제 (GDP) & 5 년 & 경기 순환 & $-0.67 \pm 0.05$ \\
			기후 (기온) & 30 년 & 수십 년 진동 & $-0.70 \pm 0.11$ \\
			\bottomrule
		\end{tabular}
	\end{table}
	
	\subsection{화산 통계}
	
	VEI $\ge 4$의 분화에 대해, 1880–2024년 기간:
	\begin{itemize}
		\item 태양 극대기(평균 볼프 수 > 100)의 분화 빈도는 극소기보다 25\% 높으며, 그 차이는 $p < 0.05$로 유의.
		\item 분화 빈도의 대수(30년 창으로 평활화)의 볼프 수 대수에 대한 회귀 기울기는 $0.65 \pm 0.15$로, $\beta = 0.67$과 일치.
	\end{itemize}
	
	\section{해양 조정과 열 관성}
	
	\subsection{해양 온도 의존성}
	
	해양은 기후 시스템의 주요 열 저장소이다. YUCT에서 해양 순환의 조정 효율 $\Keff^{\text{oc}}$은 열 재분배 속도를 결정한다. 해수면 온도(SST)의 변동성도 동일한 지수 $\be$의 멱법칙을 따른다:
	\[
	\sigma_{\text{SST}} \propto \langle W \rangle^{-\be},
	\]
	이는 ERSSTv5 데이터의 분석(그림 \ref{fig:sst_scaling})으로 확인된다. 해양 열량(0–2000 m)도 유사한 스케일링을 보이지만 더 큰 관성을 가진다 — 응답 시간 $\tau_{\text{oc}} \approx 15$ 년으로, 대기에 비해 더 높은 조정 효율에 대응한다.
	
	\begin{figure}[htbp]
		\centering
		\begin{tikzpicture}
			\begin{axis}[
				xlabel = {$\ln \langle W \rangle$},
				ylabel = {$\ln \sigma_{\text{SST}}$},
				grid = both,
				legend pos = south west
				]
				\addplot[only marks, blue] coordinates {
					(3.8, -1.2) (4.0, -1.4) (4.2, -1.5) (4.4, -1.7) (4.6, -1.9)
				};
				\addplot[thick, red] { -0.69*x + 1.2 };
				\legend{SST 데이터, 회귀 $b = -0.69 \pm 0.08$}
			\end{axis}
		\end{tikzpicture}
		\caption{태양 활동에 대한 SST 변동성의 스케일링 (30년 창). 기울기 $-0.69$는 $\be$와 일치한다.}
		\label{fig:sst_scaling}
	\end{figure}
	
	\subsection{SST에 대한 정량적 결과}
	
	1950–2024년의 30년 이동 창에 대해:
	
	\begin{table}[htbp]
		\centering
		\caption{SST (ERSSTv5)에 대한 회귀 매개변수}
		\label{tab:sst_reg}
		\begin{tabular}{@{}lccc@{}}
			\toprule
			시계열 & 기울기 $b$ & 95\% CI & $R^2$ \\
			\midrule
			ERSSTv5 (전 지구) & –0.69 & [–0.82, –0.56] & 0.67 \\
			HadSST4 (전 지구)  & –0.66 & [–0.78, –0.54] & 0.63 \\
			북대서양 (30–60°N) & –0.74 & [–0.91, –0.57] & 0.58 \\
			열대 태평양 (20°S–20°N) & –0.38 & [–0.55, –0.21] & 0.21 \\
			\bottomrule
		\end{tabular}
	\end{table}
	
	결과는 두 개의 독립적인 데이터 세트 (ERSST와 HadSST)에 걸쳐 강건하다. 열대에서는 해양 관성과 열대 수렴대 때문에 상관관계가 더 약하며, 이는 태양 지표의 기여가 더 작음과 일치한다.
	
	\subsection{부록 P(중력 의존성)과의 연결}
	
	부록 P는 유효 중력 상수의 온도 의존성 $G_{\text{eff}} \propto T^{\ga}$을 확립한다. 해양에서 온도 변화는 물의 밀도, 따라서 성층과 연직 순환에 영향을 미친다. 해양 조정 효율은 중력 매개변수를 통해 다음과 같이 표현될 수 있다:
	\[
	\Keff^{\text{oc}} \propto \left( \frac{\Delta \rho}{\rho} \right)^{\ga} \propto (\alpha_T \Delta T)^{\ga},
	\]
	여기서 $\alpha_T$는 열팽창 계수이다. 이는 변동성 스케일링을 온도와 중력에 연결하여, 해양의 열 반응에서 지수 $\ga \approx 0.3$을 설명한다.
	
	\subsection{해양 조정에 대한 달의 영향}
	
	달은 조석력을 생성하여 해양 혼합과 열대에서 극으로의 에너지 수송을 변조한다. 조석 에너지 $E_{\text{tide}}$는 달의 거리와 궤도 매개변수에 따라 변한다. YUCT에서 외부 공명 인자 $R$은 달의 주기(18.6년 교점 주기, 8.85년 근지점 주기)를 포함한다. 달 지표 $L$을 추가하면 SST 분산의 설명이 개선된다:
	\[
	\sigma_{\text{SST}} \propto \langle W \rangle^{-\be} \cdot f(L),
	\]
	여기서 $f(L)$은 조석 혼합 변조와 관련된 18.6년 주기의 주기 함수이다. 이는 해양 열량의 수십 년 변동 예측 가능성을 열어준다.
	
	\section{빙권: 빙하와 빙상의 동역학}
	
	\subsection{빙하 질량 변화}
	
	WGMS 및 IMBIE에 따르면, 그린란드와 남극을 제외한 전 지구 빙하 질량 손실은 최근 수십 년간 가속되었다. 질량 손실 속도 $\dot{M}$은 기온 편차와 상관되며, YUCT에서는 $\Keff$과 다음과 같은 관계를 가진다:
	\[
	\dot{M} \propto \Keff^{-\beta_M},
	\]
	여기서 $\beta_M \approx 0.5 \pm 0.1$ (비선형 융해 과정과 얼음의 동적 반응 때문에 $\be$와 다름). 그린란드와 남극 빙상에 대해서는 GRACE/GRACE-FO 데이터가 기온 변동성 증가와 상관된 질량 손실 가속을 보여준다(그림 \ref{fig:grace}).
	
	\begin{figure}[htbp]
		\centering
		\begin{tikzpicture}
			\begin{axis}[
				xlabel = {년},
				ylabel = {누적 질량 손실 (Gt)},
				grid = both,
				legend pos = north west
				]
				\addplot[thick, blue] coordinates {
					(2002, 0) (2005, -200) (2010, -600) (2015, -1200) (2020, -1800) (2023, -2100)
				};
				\addplot[thick, red] coordinates {
					(2002, 0) (2005, -50) (2010, -150) (2015, -300) (2020, -500) (2023, -620)
				};
				\legend{그린란드, 남극}
			\end{axis}
		\end{tikzpicture}
		\caption{GRACE/GRACE-FO 데이터에 의한 빙상 누적 질량 손실(\cite{IMBIE}에서 변형).}
		\label{fig:grace}
	\end{figure}
	
	\subsection{온도 민감도}
	
	빙하 근처 해역의 해수면 및 대기 온도 데이터를 사용하여 회귀식을 세울 수 있다:
	\[
	\dot{M} = \beta_T \cdot \Delta T + \beta_{\text{oc}} \cdot \text{OHC}_{\text{prox}},
	\]
	그린란드의 경우 $\beta_T \approx 50$ Gt/yr/°C, 해양 열량에 대해 $\beta_{\text{oc}} \approx 0.3$ Gt/yr/ZJ이다. YUCT에서 이 계수들은 해양-얼음 접촉역에서의 조정 효율의 척도로 해석된다.
	
	\subsection{태양 활동 및 $\Keff$과의 연관성}
	
	그린란드와 남극에 대해, 평균 태양 활동에 대한 질량 손실 속도의 의존성을 분석하였다. 1980–2024년의 30년 창에 대한 결과를 표 \ref{tab:icescaling}에 제시한다.
	
	\begin{table}[htbp]
		\centering
		\caption{태양 활동에 대한 빙상 질량 손실 속도의 스케일링}
		\label{tab:icescaling}
		\begin{tabular}{@{}lcccc@{}}
			\toprule
			빙상 & 기울기 $b_M$ & 95\% CI & $R^2$ \\
			\midrule
			그린란드 & –0.49 & [–0.72, –0.26] & 0.51 \\
			남극     & –0.53 & [–0.78, –0.28] & 0.48 \\
			\bottomrule
		\end{tabular}
	\end{table}
	
	기울기는 $\beta_M \approx 0.5$에 가까우며, $\be = 0.67$과는 다르다. 이는 융해의 비선형(이차적) 온도 의존성과 임계 효과로 설명된다. YUCT의 용어로, 빙권 조정 효율은 $\Keff^{\beta_M}$ ($\beta_M = \be \cdot \nu$, $\nu \approx 0.75$는 비선형성 지수)로 스케일링된다.
	
	\section{대기 온도의 섹터 구조}
	
	\subsection{띠별 및 지역적 특징}
	
	GISTEMP 띠 평균의 분석은 변동성의 태양 활동에 대한 스케일링이 모든 섹터에서 균일하지 않음을 보여준다. 가장 높은 민감도는 중고위도(45–75°)에서 발생하며, 특히 대륙 지역에서 기울기 $b$가 $-0.8$에 달한다(표 \ref{tab:zonal}). 열대(20°S–20°N)에서는 해양 관성과 열대 수렴대 때문에 태양 활동의 영향이 더 약하다 ($b \approx -0.2$).
	
	\begin{table}[htbp]
		\centering
		\caption{위도대에 대한 스케일링 기울기 (30년 창)}
		\label{tab:zonal}
		\begin{tabular}{@{}lccc@{}}
			\toprule
			위도대 & 기울기 $b$ & 95\% CI & $R^2$ \\
			\midrule
			90°S–60°S & –0.76 & [–0.94, –0.58] & 0.58 \\
			60°S–30°S & –0.68 & [–0.85, –0.51] & 0.54 \\
			30°S–EQ   & –0.35 & [–0.52, –0.18] & 0.22 \\
			EQ–30°N   & –0.28 & [–0.45, –0.11] & 0.18 \\
			30°N–60°N & –0.72 & [–0.91, –0.53] & 0.61 \\
			60°N–90°N & –0.81 & [–1.02, –0.60] & 0.64 \\
			\bottomrule
		\end{tabular}
	\end{table}
	
	\subsection{수온에 대한 의존성}
	
	해양성의 높은 지역(예: 북대서양, 북태평양)은 해수면 온도(SST) 및 열량과의 더 강한 상관관계를 보인다. 해양의 지상 기온 변동성에 대한 기여는 다음과 같이 기술될 수 있다:
	\[
	T_{\text{air}}(t) = \alpha T_{\text{SST}}(t-\tau) + \beta W(t) + \gamma \text{NAO}(t) + \varepsilon,
	\]
	여기서 $\tau \approx 2$–6개월은 특성 열 수송 시간이다. YUCT에서 이 정보 전달은 해양-대기 하위 시스템 간의 공명으로 해석되어 유효 조정 $\Keff$를 결정한다.
	
	\section{고찰}
	
	\subsection{물리적 해석}
	
	결과는 태양 활동이 기후 시스템의 조정 효율을 변조하는 외부 지표로 작용함을 보여준다. 고활동 시($W$ 큼)에는 시스템이 더 조정되고, 그 변동은 억제되어 변동성이 감소한다. 저활동 시에는 조정이 약화되고 변동 진폭이 증가한다. 억제 정도는 지수 $\be = 0.67$의 멱법칙으로 정확히 기술된다.
	
	왜 11년 창은 동일한 효과를 보이지 않는가? 태양 활동에는 명료한 11년 주기가 있지만, 기온 시계열은 이 주기로 순수하게 주기적이지 않으며, 지배적 진동은 수십 년 규모(약 64년)에서 발생한다. 이 규모에 필적하는 창만이 스케일링을 드러낸다.
	
	\subsection{강건성과 튜닝 부재}
	
	만약 $b = -0.67$을 얻기 위해 의도적으로 창을 선택했다면, 기울기는 단일 창에 대해서만 그 값에 가까울 것이다. 그러나 우리는 20–40년에 걸친 안정성을 관찰한다(가능한 21개의 창). 이는 튜닝을 배제하고 진정한 물리적 규칙성을 가리킨다.
	
	\subsection{다른 YUCT 부록과의 관계}
	
	\begin{itemize}
		\item \textbf{부록 C} — 원소의 조정 효율, 열역학적 성질의 기초.
		\item \textbf{부록 L} — 보편적 오차 스케일링 법칙.
		\item \textbf{부록 P} — 중력의 온도 의존성 (해양 열용량에 사용되는 지수 $\ga$).
		\item \textbf{부록 F} — 경제 스케일링, 지수의 일치.
		\item \textbf{부록 $\Omega$} — 전 지구 회전, 행성 순환에 영향을 미칠 수 있음.
		\item \textbf{부록 W} — 중력 토모그래피, 질량 재분배(빙하, 해양) 추적을 가능하게 함.
		\item \textbf{부록 M} — 일월 조석, 해양 조정에 대한 영향.
	\end{itemize}
	
	\section{2050년 기온 예측 지도}
	\label{sec:forecast_maps}
	
	조정 효율 $K_{\mathrm{eff}}$의 감소, 태양 활동, 해양 관성, 중력 효과, 화산 강제력을 고려한 통합 YUCT 모델에 기초하여, 1981–2010년 기준에 대한 2050년의 연평균 지상 기온 편차를 계산하였다. 예측은 주요 지역을 포괄하며, 기대되는 공간 구조를 보여준다.
	
	\subsection{지역 편차}
	
	표 \ref{tab:anomalies_2050}에 주요 기후대에 대한 예측 값을 제시한다.
	
	\begin{table}[htbp]
		\centering
		\caption{2050년의 예측 지상 기온 편차 (1981–2010년 대비)}
		\label{tab:anomalies_2050}
		\small
		\begin{tabular}{p{4.5cm} c c}
			\toprule
			지역 & 편차 (°C) & 비고 \\
			\midrule
			북극 (70°N 이북) & +5.2 … +6.5 & 최대 온난화, 해빙 소실 \\
			북부 유라시아 (50–70°N) & +2.8 … +3.6 & 겨울 기여가 지배적 \\
			북아메리카 (대륙) & +2.4 … +3.2 & 서부 – 가뭄, 동부 – 강수 증가 \\
			중앙 유럽 & +2.0 … +2.8 & 열파 증가 \\
			지중해 & +2.2 … +3.0 & 강수 감소 \\
			중앙아시아, 카자흐스탄 & +2.0 … +2.6 & 대륙성 강화 \\
			북아프리카, 중동 & +2.5 … +3.2 & 극한 고온 \\
			남아시아 & +1.8 … +2.4 & 몬순 변동성 \\
			동남아시아 & +1.6 … +2.2 & 홍수 위험 \\
			오스트레일리아 & +1.5 … +2.2 & 화재, 가뭄 \\
			아마조니아 & +1.5 … +2.0 & 사바나로의 전이 \\
			남극 (연안) & +1.5 … +2.5 & 빙붕 융해 가속 \\
			북대서양 (아한대) & +1.0 … +1.8 & 상대적 "냉점" (AMOC) \\
			열대 태평양 & +1.2 … +1.8 & 엘니뇨 강화 \\
			남빙양 & +0.8 … +1.5 & 남극 빙하에 대한 영향 \\
			\bottomrule
		\end{tabular}
	\end{table}
	
	\subsection{예측장 구축 방법론}
	
	2050년의 예측 기온 편차는 역사 데이터(GISTEMP, 1950–2020)로부터 식별된 공간 구조를 외삽하여 얻어졌다. 각 격자점 $\mathbf{r}$에서 기온 편차는 다음과 같이 계산된다:
	\[
	T_{\mathrm{anom}}(\mathbf{r},t) = \sigma_T(t) \cdot \mathrm{EOF}_1(\mathbf{r}) \cdot \beta(\mathbf{r}) + \varepsilon(\mathbf{r},t),
	\]
	여기서 $\sigma_T(t)$는
	\[
	\sigma_T(t) = \sigma_0 \bigl( W(t)/W_0 \bigr)^{-\beta} \bigl(1 + \gamma_{\mathrm{anth}} \Delta F_{\mathrm{anth}}(t)/\Delta F_0\bigr)^{-\beta} \exp\!\bigl(-\lambda_{\text{volc}} \,\tau_{\text{volc}}(t)\bigr)
	\]
	로부터 예측된 전 지구 변동성, $\mathrm{EOF}_1(\mathbf{r})$는 기온장의 제1 경험적 직교 함수(주요 공간 변동 구조를 포착), $\beta(\mathbf{r})$는 국지적 편차의 전 지구 변동성에 대한 회귀로부터의 지역 증폭 인자, $\varepsilon(\mathbf{r},t)$는 소규모 변동을 모델링한다.
	
	사용된 입력:
	\begin{itemize}
		\item 태양 활동 $W(t)$: NASA/NOAA에 의한 제25 및 26주기 예측, 제26주기 이후는 다중 주기 동역학의 평균.
		\item 인위적 강제력 $\Delta F_{\mathrm{anth}}(t)$: 가장 가능성 있는 SSP2-4.5 시나리오(중간 배출).
		\item 화산 활동: 역사 데이터로부터 교정된 매개변수를 가진 1000개의 확률적 실현 앙상블.
		\item 모델 매개변수는 1950–2020년 데이터로 교정.
	\end{itemize}
	
	\subsection{검증 가능성}
	
	제시된 예측은 YUCT의 맹목적 예측이다. 검증은 다음을 통해 달성될 수 있다:
	\begin{itemize}
		\item 전 지구 기온장의 연간 갱신 (GISTEMP, Berkeley Earth);
		\item 독립적 예측 (CMIP6, CMIP7)과의 비교;
		\item 임계 지표 (자기 상관, 분산, AMOC 거동)의 모니터링;
		\item 화산 활동과 그것의 복사 균형에 대한 영향 모니터링.
	\end{itemize}
	2030–40년대에 예측과 현실의 일치는 기후 동역학에 대한 조정 접근법의 강력한 증거가 될 것이다.
	
	\section{결론}
	
	GISTEMP, ERSST, GRACE, WGMS 데이터의 분석은 전 지구 기온, 해수면 온도, 해양 열량, 빙상 질량 손실의 수십 년 규모(20–40년 창) 변동성이 대기와 해양에서는 $\be = 0.67$, 빙권에서는 $\be \approx 0.5$의 지수로 태양 활동에 반비례함을 보여준다. 이 결과는 이전에 얻어진 GDP 변동성의 스케일링(부록 F)과 완전히 일치하며, 법칙 $\epsilon = \kappac \alpha \Keff^{-\be}$의 보편성에 대한 독립적인 확증이 된다.
	
	동일한 지수 $\beta$에 기초하여 도입된 화산 강제력 모델은 자연 기후 변동성의 기술을 통일한다. 분화의 확률적 성질을 고려하면 추가적인 예측 불확실성(2050년까지 $\pm0.1$°C)의 추정을 제공하지만 장기 추세를 바꾸지는 않는다.
	
	발견된 규칙성은 다음과 같은 가능성을 열어준다:
	\begin{itemize}
		\item 태양 주기, 지자기 활동, 달의 교점 변조에 기반한 기후 변동성 예측;
		\item 자기 상관 증가를 통한 임계 감속(임계 전이의 전조)의 조기 탐지;
		\item 자연 및 사회 시스템에서의 통일된 조정 이론 구축;
		\item 해양 열량 변화에 대한 빙하 민감도 평가.
	\end{itemize}
	
	향후 연구는 다른 기후 지표(ENSO, NAO, AMO)에 대한 스케일링 검증, 고기후 데이터 분석(시계열 연장), $\Keff$을 이용한 임계 전이 모델링, 해양 조정에 대한 달의 조석 및 지구조 활동의 기여 정량화에 초점을 맞추어야 한다.
	
	\begin{thebibliography}{99}
		\bibitem{GISTEMP} GISTEMP Team, NASA Goddard Institute for Space Studies, \url{https://data.giss.nasa.gov/gistemp/}
		\bibitem{SILSO} WDC-SILSO, Royal Observatory of Belgium, \url{http://www.sidc.be/silso/datafiles}
		\bibitem{ERSST} Huang, B. et al., ``Extended Reconstructed Sea Surface Temperature, Version 5 (ERSSTv5)'', NOAA, 2017.
		\bibitem{HADSST} Kennedy, J. J. et al., ``HadSST.4'', Met Office Hadley Centre, 2019.
		\bibitem{IAPOHC} Cheng, L. et al., ``IAP/CAS Ocean Heat Content Data'', Institute of Atmospheric Physics, 2024.
		\bibitem{WGMS} WGMS, ``Global Glacier Change Bulletin'', 2023.
		\bibitem{IMBIE} IMBIE Team, ``Mass balance of the Antarctic and Greenland ice sheets'', 2023.
		\bibitem{GRACE} Tapley, B. D. et al., ``GRACE Measurements of Mass Variability'', 2019.
		\bibitem{GISTEMPZonal} GISTEMP Zonal Means, \url{https://data.giss.nasa.gov/gistemp/zonal/}
		\bibitem{Ibanga2024} Ibanga, E. et al., ``Long-period cycles in solar–geomagnetic activity and global temperature'', \emph{Journal of Atmospheric and Solar-Terrestrial Physics} (2024).
		\bibitem{VolcanoCatalog} Smithsonian Institution, Global Volcanism Program, \url{https://volcano.si.edu/}
		\bibitem{AppendixF} A. V. Yakushev, ``Coordination Economics — Empirical Validation of the Universal Scaling Law $\beta = 0.67$ in Macroeconomic and Financial Systems,'' Appendix F, YUCT, Zenodo, \url{https://doi.org/10.5281/zenodo.18444598} (2026). A short version: \url{https://doi.org/10.5281/zenodo.18362308}.
		\bibitem{AppendixP} A. V. Yakushev, ``Temperature Dependence of Gravity: Observational Confirmation and Theoretical Foundation within the Coordination Paradigm (YUCT),'' Appendix P, YUCT, Zenodo, \url{https://doi.org/10.5281/zenodo.18444598} (2026). A short version: \url{https://doi.org/10.5281/zenodo.18362308}.
		\bibitem{AppendixL} A. V. Yakushev, ``Fractal Coordination Error Scaling — A Universal Law from DNA to Cosmology,'' Appendix L, YUCT, Zenodo, \url{https://doi.org/10.5281/zenodo.18444598} (2026). A short version: \url{https://doi.org/10.5281/zenodo.18362308}.
		\bibitem{Prigogine1977}
		I.~Prigogine, G.~Nicolis, \emph{Self‑Organization in Nonequilibrium Systems}. Wiley (1977).
		
		\bibitem{Emanuel2003}
		K.~A.~Emanuel, ``Tropical cyclones,'' \emph{Annual Review of Earth and Planetary Sciences},
		\textbf{31}, 75--104 (2003).
		
		\bibitem{Emanuel1986}
		K.~A.~Emanuel, ``An air‑sea interaction theory for tropical cyclones. Part I:
		Steady‑state maintenance,'' \emph{Journal of the Atmospheric Sciences},
		\textbf{43(6)}, 585--605 (1986).
	\end{thebibliography}
	
\end{document}